\documentclass{article}

\PassOptionsToPackage{numbers, compress}{natbib}

\usepackage[main]{neurips_2026}

% Encoding
\usepackage[utf8]{inputenc}
\usepackage[T1]{fontenc}

% Core math
\usepackage{amsmath, amssymb, amsfonts}
\usepackage{mathtools}
\usepackage{bm}

% Figures and tables
\usepackage{graphicx}
\usepackage{booktabs}
\usepackage{tabularx}
\usepackage{multirow}
\usepackage{makecell}
\usepackage{colortbl}
\usepackage{siunitx}
\usepackage{rotating}
\usepackage{wrapfig}

% Formatting
\usepackage{microtype}
\usepackage{xcolor}
\usepackage{nicefrac}
% \usepackage{refcheck}

% Links (load once, near the end)
\usepackage[bookmarks=true]{hyperref}

% Algorithms and code
\usepackage[ruled]{algorithm2e}
\usepackage{listings}

% Subfigures
\usepackage[caption=false,font=footnotesize]{subfig}

% Misc
\usepackage{multicol}
\usepackage{dsfont}
\usepackage{tabularray}

% Note. For the workshop paper template, both \title{} and \workshoptitle{} are required, with the former indicating the paper title shown in the title and the latter indicating the workshop title displayed in the footnote. 
\title{Learning to Spray in an Uncertain and Windy Environment}


\author{%
  David S.~Hippocampus\thanks{Use footnote for providing further information
    about author (webpage, alternative address)---\emph{not} for acknowledging
    funding agencies.} \\
  Department of Computer Science\\
  Cranberry-Lemon University\\
  Pittsburgh, PA 15213 \\
  \texttt{hippo@cs.cranberry-lemon.edu} \\
}

% Add macros
% Formatting and General
\newcommand{\etal}{\textit{et al. }}
\newcommand{\solution}{\textbf{Solution:}}
\newcommand{\qed}{\hfill \ensuremath{\Box}}
\definecolor{darkgreen}{rgb}{0.0, 0.5, 0.0} 

% Math Sets and Symbols
\newcommand{\reals}{\ensuremath{\mathbb{R}}}
\newcommand{\naturals}{\ensuremath{\mathbb{N}}}
\newcommand{\integers}{\ensuremath{\mathbb{Z}}}
\newcommand{\expectation}{\ensuremath{\mathbb{E}}}
\newcommand{\D}{\mathrm{d}}
\newcommand{\reg}{\ensuremath{\emph{R}}}

% Spaces and Sets
\newcommand{\config}{\mathcal{Q}}
\newcommand{\wspace}{\mathcal{W}}
\newcommand{\obst}{\mathcal{O}}
\newcommand{\goals}{\mathcal{G}}
\newcommand{\states}{\mathcal{S}}
\newcommand{\trans}{\ensuremath{\mathcal{T}}}

% Rewards
\newcommand{\rterm}{R_{\mathrm{term}}}
\newcommand{\rspray}{R_{\mathrm{spr}}}
\newcommand{\rpath}{R_{\mathrm{path}}}
\newcommand{\rvisit}{R_{\mathrm{visit}}}

% Theorem Environments
\newtheorem{problem}{Problem}
\newtheorem{remark}{Remark}
\newtheorem{definition}{Definition}
\newtheorem{proposition}{Proposition}
\newtheorem{lemma}{Lemma}
\newtheorem{theorem}{Theorem}

% Paired Delimiters
\DeclarePairedDelimiter\ceil{\lceil}{\rceil}
\DeclarePairedDelimiter\floor{\lfloor}{\rfloor}

% Operators
\DeclareMathOperator*{\argmax}{arg\,max}
\DeclareMathOperator{\simulate}{SIM}
\DeclareMathOperator{\init}{init}
\DeclareMathOperator{\free}{free}
\DeclareMathOperator{\new}{new}
\DeclareMathOperator{\pos}{POS}
\DeclareMathOperator{\goalf}{GS}
\DeclareMathOperator{\visited}{\mathbb{V}}

% Fallback definitions for general-formulation macros (no-op if already defined in macros.tex)
\providecommand{\obst}{\mathcal{O}}
\providecommand{\free}{\mathrm{free}}
\providecommand{\obs}{\mathrm{obs}}
\providecommand{\col}{\mathrm{col}}

\begin{document}

\maketitle

\begin{abstract}
The deployment of robotic systems for selective spraying offers a sustainable alternative to blanket spraying in agriculture. However, existing precision spraying systems typically assume deterministic kinematics and perfect observations, ignoring the inherent stochasticity of agricultural fields, such as unpredictable wind patterns, sensor noise, and variable spray-dispersion dynamics. To bridge this gap, we introduce a novel continuous-control environment for multi-robot selective spraying that explicitly models real-world uncertainties: stochastic wind, actuation noise, observation noise, and probabilistic spray decay. We formulate the task as an optimization problem—coordinating robots to disinfect targets under uncertainty while minimizing path lengths and avoiding collisions—and prove its equivalence to a Partially Observable Markov Decision Process (POMDP). Our environment is compatible with standard deep Reinforcement Learning (RL) algorithms and serves as a rigorous benchmark for robust agricultural coordination, which can be modified or extended for other similar applications with real-world uncertainties. We evaluate several state-of-the-art deep RL methods across multiple environment variants, demonstrating the feasibility of learning effective policies under realistic uncertainty.
\end{abstract}

\section{Introduction}
% \label{sec:intro}

Effective pest and weed management is critical to agricultural productivity, yet traditional blanket spraying causes environmental degradation and health risks \citep{zhou2025ec, mengistie2025, shekhar2024ai, jugulam2016}. Selective spraying directly mitigates these issues; robotic spot-sprayers and UAVs have demonstrated chemical reductions of 35\% to 60\% while maintaining efficacy \citep{saleh2025spotspray, allmendinger2022cnn}. Deploying single and multi-robot systems for precision spraying therefore offers a highly promising path toward sustainable farming \citep{ratan2021, ratan2017, baltazar2021smarter, ju2022cea, saleh2025spotspray, Bzhikhatlov2023}. However, real-world deployment is fundamentally bottlenecked by \emph{environmental uncertainties}.  Wind disturbances alter robot trajectories and deflect spray droplets away from intended targets; actuator imprecision causes commanded thrusts to deviate from their nominal values; chemical deposition is inherently uncertain due to turbulence and canopy interaction; and onboard sensors provide only noisy estimates of robot states and infection levels. Ignoring these factors during environment design produces policies that are brittle when transferred to the field---a well-documented failure mode in applied reinforcement learning (RL) \citep{dulacArnold2021}. While much research has been done on the improvement of RL algorithms, the design and development of RL environments is equally important and comes with unique challenges \citep{reda2020learning, dulac2020empirical, jakub2023as, henderson2018deep, sutton2018reinforcement}.

Many existing approaches to multi-robot precision spraying formulate the task as a combinatorial or mixed-integer program and compute open-loop plans under deterministic assumptions \citep{ratan2021, ratan2017, xu2023multiBases, xu2022jointVehicle, li2022optimal, huang2024allocation}. While these methods are principled, they require the field geometry and infection locations to be fixed, scale poorly to high-dimensional continuous spaces, and provide no mechanism for adapting to disturbances encountered during execution. On the learning-based side, RL has recently been applied to single-robot precision spraying \citep{hao2022rl, khosravi2025hamppo}, with \citet{khosravi2025hamppo} demonstrating a hierarchical PPO policy that jointly learns navigation and spray application for biotic stress management in row-crop fields. However, these works use deterministic simulation environments: no wind model is included, actuation noise and spray deposition uncertainty are absent, and the framework is designed for a single robot. Broader surveys of RL in agricultural robotics \citep{digitaltwinRL2024, uavRL2025} confirm that no existing work has applied RL to multi-robot selective spraying in a stochastic, noisy environment.

This paper fills that gap by formalizing an efficient, customizable environment for multi-robot selective spraying. We formulate the task as an optimization problem, develop the equivalent POMDP, and prove that the optimal policy of the POMDP solves the optimization problem. Our specific contributions are:
\begin{enumerate}
    \item \textbf{Stochastic multi-robot spraying environment.} We present a novel environment for cooperative multi-robot selective spraying that incorporates stochastic wind dynamics, Gaussian spray-deposition, actuation noise, observation noise, and per-episode domain randomization.

    \item \textbf{Optimization-to-POMDP equivalence.} We formally prove that the POMDP's optimal policy solves the constrained optimization objective, providing theoretical grounding absent in prior works \citep{ratan2017, ratan2021, khosravi2025hamppo}.

    \item \textbf{Comprehensive benchmarking.} We evaluate six state-of-the-art deep RL algorithms---A2C \citep{mnih2016a2c}, ARS \citep{mania2018ars}, TRPO \citep{schulman2015trust}, PPO \citep{schulman2017ppo}, TQC \citep{kuznetsov2020tqc}, and CrossQ \citep{bhatt2024crossq}---across ten environment variants, demonstrating that standard algorithms generalize to the stochastic multi-robot spraying setting without modification.

    \item \textbf{High-fidelity 3-D simulation.} We simulate the environment on an agricultural field with herbicide-spraying UAVs using the CoppeliaSim robot simulator \citep{coppeliaSim2013} with Pybullet \citep{coumans2021pybullet} physics engine, providing a bridge between simulation and physically realistic deployment.
\end{enumerate}

\section{Related Work}
\label{sec:recent}

\subsection{Precision Spraying in Agriculture}
Selective spraying significantly reduces chemical waste compared to blanket application \citep{UTSTUMO2018cea, baltazar2021smartear, ju2022cea}. Single-robot systems have evolved from GPS-guided maps \citep{GonzalezdeSantos2016, allmendinger2022cnn} to real-time onboard deep learning detection, achieving up to 60\% agrochemical reduction \citep{nasir2023tobacco, saleh2025spotspray, chen2024fourwheel}. However, these rely on fixed, rule-based controllers that lack dynamic adaptability. Multi-robot coordination typically utilizes optimization-based methods under discrete and/or deterministic assumptions \citep{ratan2021, ratan2017, xu2023multiBases, xu2022jointVehicle, li2022optimal, huang2024allocation, ju2025ugv}. These open-loop approaches perform well in structured settings but lack the closed-loop adaptability needed to handle stochastic disturbances such as wind, actuation noise, or uncertain spray deposition.

\subsection{Task Allocation and Planning}
PDDL \citep{fox2003pddl21} and PDDL+ \citep{fox2006pddlplus} have been widely used for multi-robot planning with temporal constraints, with planners such as POPF \citep{coles2010popf} and OPTIC \citep{benton2012optic} handling continuous numeric effects. ScottyActivity \citep{fernandez2018scottyactivity} interleaves heuristic search with convex optimization for coordinated teams, but explicitly assumes first-order linear dynamics and no runtime uncertainty. Conflict-Based Search (CBS) \citep{sharon2015cbs} and its extensions \citep{li2019cbscontinuous} provide optimal or bounded-suboptimal multi-agent path finding on discretized graphs with deterministic edge costs. Sampling-based methods with robot dynamics \citep{le2017icaps, le2019ral, le2021icra, bui2022iros} construct deterministic motion trees over precisely modeled dynamics and do not model stochastic disturbances or partial observability. All of these methods produce plans susceptible to compounding errors under the continuous stochastic perturbations present in real agricultural environments. By optimizing expected cumulative reward under the full distribution of dynamics, RL offers a substantially more robust alternative for stochastic and noisy environments.

\subsection{Reinforcement Learning}
Reinforcement learning (RL) has been increasingly applied to agricultural robotics for navigation, monitoring, and decision-making tasks \citep{hao2022rl, digitaltwinRL2024}. For spraying, \citet{hao2022rl} proposed an adaptive spraying decision system in which a UAV learns spraying parameters based on crop conditions, achieving reductions in pesticide usage. However, this work focuses on parameter optimization rather than joint navigation and spraying, and does not consider multi-robot coordination or environmental stochasticity. The most closely related learning-based approach is HAM-PPO \citep{khosravi2025hamppo}, which uses a hierarchical PPO framework to jointly learn navigation and spraying decisions for a single robot. However, the simulation environment is simplified and does not explicitly model wind disturbances, actuation noise, or stochastic spray deposition. Existing surveys of RL in agriculture and UAV systems \citep{digitaltwinRL2024, uavRL2025} indicate that RL-based spraying applications are largely limited to single-agent settings with simplified or deterministic environments.

In broader robotics, RL has achieved strong performance in multi-robot navigation and collision avoidance \citep{chen2017icra, chen2021iral, everett2018iros, fan2020ijrr}. These works focus on decentralized control and safe navigation, often in environments with static or dynamic obstacles, but do not address spraying tasks or consider such sources of environmental uncertainty. Recent advances in multi-agent reinforcement learning (MARL) provide scalable frameworks for decentralized coordination under uncertainty. Value decomposition methods such as VDN \citep{sunehag2018vdn} and QMIX \citep{rashid2018qmix, rashid2020qmix} enable tractable learning by factorizing joint value functions, while actor–critic approaches such as MADDPG \citep{lowe2017maddpg} and MAPPO \citep{yu2022mappo} achieve strong empirical performance in cooperative settings. These methods have been successfully applied to large-scale coordination problems, but typically assume simplified dynamics or synthetic benchmarks, and lack formal correctness guarantees that limit their direct applicability to physically grounded agricultural environments with stochastic disturbances. 

\subsection{Environment Design}

The importance of carefully designed environments is well recognized \citep{reda2020learning, cobbe2020icml, dulac2020empirical, sutton2018reinforcement, henderson2018deep}. Benchmarking platforms such as Google Research Football \citep{kurach2020google} and CityFlow \citep{zhang2019cityflow} provide environments for standardized tasks, but neither includes formal descriptions or guarantees. Physics-based simulators such as MuJoCo \citep{todorov2012mujoco}, PyBullet \citep{coumans2021pybullet}, and Isaac Gym \citep{makoviychuk2021isaac} enable more realistic modeling of robot dynamics and contact interactions, and have been widely used in robotics control and sim-to-real transfer studies \citep{tassa2018deepmind, akkaya2019solving, tobin2017domain}. This motivates us to develop a custom, Gymnasium-compatible \citep{towers2024gymnasium} 3-D real-time environment using the PyBullet \citep{coumans2021pybullet} physics engine within the CoppeliaSim \citep{coppeliaSim2013} robotics simulator. Unlike existing methods, our cooperative multi-robot formulation explicitly incorporates the multiple sources of physical uncertainty inherent to real agricultural fields—including stochastic wind dynamics, noisy actuation, and uncertain spray deposition—into a continuous 3-D control loop. We believe that our proposed environment, along with its solution using state-of-the-art deep RL methods, will give researchers a new direction to be able to deploy formally verified RL-based systems that autonomously manage not just selective spraying, but a wide range of coordination tasks, leading to highly efficient and sustainable robotic systems.

\section{Preliminaries}
% \label{sec:preli}
Let $\reals, \reals_{\geq 0}, \naturals$ be the set of real, non-negative real, and natural numbers, respectively. Let $\{x_i\}_{i=1}^n$ denote an ordered sequence or tuple of elements $x_1, x_2, \ldots, x_n$. For convenience, the $L^2$ or Euclidean norm $\|.\|_2$ is denoted simply as $\|.\|$. Let $\mathbf{I}_n$ denote the identity matrix of dimension $n \times n$. Let $\mathbf{1}_n$ and $\mathbf{0}_n$ denote $n$-dimensional all-ones and all-zeros vectors, respectively. Given two real numbers $x, y \in \reals$, we define the intervals $[x,y] = \{ i \in \reals \mid x \leq i \leq y \}, [x,y) = \{ i \in \reals \mid x \leq i < y \}$ and $(x,y] = \{ i \in \reals \mid x < i \leq y \}$. Let $\mathds{1}_C$ be a boolean variable where $\mathds{1}_C=1$ if condition $C$ is true and $0$ otherwise. For a set of samples $S$, let $\Delta(S)$ denote the set of probability distributions over $S$. For a distribution $\mathcal{X}$ with mean $\mu \in \reals$ and variance $\sigma^2 \in \reals_{\geq 0}$, a sample $x$ drawn from this distribution is denoted as, $x \sim \mathcal{X}(\mu, \sigma^2)$. Let $\mathcal{N}(\mu, \sigma^2)$ denote the univariate normal (Gaussian) distribution. Let $\mathcal{N}(\bm{\mu}, \mathbf{\Sigma})$ denote the multivariate Gaussian distribution with mean vector $\bm{\mu} \in \reals^n$ and covariance matrix $\mathbf{\Sigma} \in \reals^{n \times n}$, where $\mathbf{\Sigma}$ is symmetric positive semi-definite \citep{bishop2006pattern}. The isotropic case $\mathbf{\Sigma} = \sigma^2 \mathbf{I}_n$ yields $n$ independent components each with variance $\sigma^2$. For a continuous random variable $Y$ with probability density function $f$, the expectation is defined as $\expectation[Y] = \int_{-\infty}^{\infty} y f(y)dy$. In RL, a policy $\pi: S \to \Delta (A)$ is a function that maps a state $s \in S$ to a probability distribution over the set of actions $A$. The action-value or Q-value of a policy $\pi$ at state $s$ and action $a$ is defined as the expected cumulative reward from following the policy after taking the action $a$ starting from $s$ and written as, $Q^{\pi}(s,a) = \expectation_{\pi}[\sum_{\tau\geq0}\gamma^\tau r_\tau | s_0=s, a_0=a]$ where $r_\tau \in \reals$ is the reward at time $\tau$ and $\gamma \in [0,1]$ is the discount factor. Following the notations and definitions from \citet{choset2005}, we define the multi-robot selective spraying problem in the next section.

\section{Problem Formulation}
% \label{sec:prob}

Consider a group of $n \in \naturals $ mobile robots $\mathcal{M} = \{ \mathcal{M}_1, \ldots, \mathcal{M}_n \}$ operating in a bounded, $d$-dimensional workspace $\wspace \subseteq \reals^d$. From \citet{le2021icra}, each robot $\mathcal{M}_i$ is modeled as $
\mathcal{M}_i = \{ \mathcal{A}_i, \config_i, \states_i, s_i^{\init}, \mathcal{U}_i, f_i \}$ in terms of its shape $\mathcal{A}_i$, configuration space $\config_i \subseteq \reals^{d+1}$, state space $\states_i = \config_i \times \mathcal{V}_i \subseteq \reals^{2d+1}$ where $\mathcal{V}_i \subseteq \reals^d$ is the motion space, initial state $s_i^{\init} \in \states_i$, control space $\mathcal{U}_i \subseteq \reals^{d+1}$, and dynamics $f_i$ (subject to stochastic disturbances). We model the robots as omnidirectional unmanned aerial vehicles (UAVs) operating in an uncertain, noisy environment. Each robot is equipped with a spraying mechanism and a finite spray capacity $c_i \in \reals_{\geq 0}$. Let $\{ c_i^{\init} \}_{i=1}^n$ be the initial spraying capacities of the $n$ robots. The configuration $q_i = \langle p_i, c_i \rangle \in \config_i$ denotes the position $p_i \in \reals^d$ and remaining spray capacity $c_i$ of robot $i$. The state $s_i = \langle p_i, c_i, v_i \rangle \in \states_i$ denotes the operational state where $v_i \in \mathcal{V}_i$ is the velocity. The control input $u_i = \langle \mu_i, \beta_i \rangle \in \mathcal{U}_i$ specifies the normalized $d$-dimensional thrust components $\mu_i \in [-1, 1]^d$ for changing the robot velocity, and the normalized spray signal $\beta_i \in [0, 1]$. The environment experiences a stochastic wind with magnitude $v_w \sim \mathcal{N}(\bar{v}_w, \epsilon_w^2)$ and direction $\theta_w \sim \mathcal{N}(\bar{\theta}_w, \epsilon_\theta^2)$, yielding a wind vector $\mathbf{w} = \langle v_w \cos\theta_w, v_w \sin\theta_w \rangle $ for $d=2$ (and $\mathbf{w} = \langle v_w \cos\theta_w, v_w \sin\theta_w, 0 \rangle $ for $d=3$). Here $\bar{v}_w$ and $\bar{\theta}_w$ are the nominal wind magnitude and direction, and $\epsilon_w^2$ and $\epsilon_\theta^2$ are the corresponding noise variances, respectively. The thrust components are scaled by a thrust power coefficient $\kappa > 0$. The actions and positions are perturbed by an action noise vector $\epsilon_a \sim \mathcal{N}(\mathbf{0}_d, \sigma_a^2 \mathbf{I}_d)$ and a position noise vector $\epsilon_p \sim \mathcal{N}(\mathbf{0}_d, \sigma_p^2 \mathbf{I}_d)$ where $\sigma_a^2$ and $\sigma_p^2$ are the action and position noise variances, respectively. Therefore, the velocity and position of robot $i$ evolve after timestep $dt$ as:
\begin{align}
    v_i(t+dt) &= \min\Big( \max\Big( v_i(t) + \tfrac{\kappa \mu_i(t) + \epsilon_a}{W_i} dt, v_{\min}\Big), v_{\max} \Big), \label{eq:vel_update} \\
    p_i(t+dt) &= p_i(t) + (v_i(t) + \mathbf{w}(t)) dt + \epsilon_p, \label{eq:pos_update}
\end{align}
where $W_i>0$ is the mass of robot $i$, $\min$ and $\max$ are applied element-wise, $v_{\min} < v_{\max}$ and $ v_{\min}, v_{\max}\in \reals$. For selective spraying, the workspace $\wspace$ represents an agricultural field containing $m \in \naturals$ spatially distributed infection sites $\mathcal{I} = \{ \mathcal{I}_1, \ldots, \mathcal{I}_m \}$. Each infected site $\mathcal{I}_j$ is associated with a fixed location $p_j \in \wspace$ and a non-negative infection intensity $\ell_j(t) \in \reals_{\geq 0}$ that evolves over time. Let $ \ell^{\init} = \{ \ell_j^{\init} \}_{j=1}^m$ be the initial infection intensities of the $m$ infected sites and $L^{\init} = \sum_{j=1}^m\ \ell_j^{\init}$ be the total initial infections. Let $\ell(t) = \{ \ell_j(t) \}_{j=1}^m$ be the sequence of infection levels at time $t$. The effect of taking spray action $\beta_i$ by robot $i$ on infection site $\mathcal{I}_j$ at time $t$ is modeled using a Gaussian falloff:
\begin{equation}
\Delta \ell_j^i(t) = \min \Big( \max\Big(0,  \Big( \beta_i(t) \cdot \exp\Big(-\frac{\| p_i(t) - p_j \|^2}{2\beta_s^2}\Big) + \epsilon_s \Big) dt  \Big), c_i(t), \ell_j(t) \Big),  \label{eq:spray}
\end{equation}
where $\epsilon_s \sim \mathcal{N}(0, \sigma_s^2)$ is the spray noise with variance $\sigma_s^2$, and $\beta_s \in \reals_{\geq 0}$ is a spray dispersion parameter. Finally, each infection level $\ell_j$ and robot capacity $c_i$ evolve after timestep $dt$ as:
\begin{align}
\ell_j(t + dt) &= \max \Big( 0, \ell_j(t) - \sum_{i=1}^n \Delta \ell_j^i(t) \Big), \label{eq:iloc_update} \\
c_i(t + dt) &= \max \Big( 0, c_i(t) - k_{\beta} \beta_i(t)dt \Big), \label{eq:cap_update} %- \sum_{j=1}^m \Delta \ell_j^i(t). 
\end{align}
where $k_{\beta} \in \reals_{\geq 0}$ is a spray coefficient. The stochastic dynamics from Eq. \eqref{eq:vel_update}--\eqref{eq:pos_update}, and infection/capacity updates from Eq. \eqref{eq:iloc_update}--\eqref{eq:cap_update} correspond to a forward Euler discretization of a continuous-time process \citep{feedback2020}. From a state $s_i(t) = \langle p_i, c_i, v_i \rangle$ at time $t$, robot $i$ can reach to a new state by applying a control $u_i(t) = \langle \mu_i, \beta_i \rangle $ for one step $dt$:
\begin{equation}
    s_i(t+dt) \gets \simulate(s_i(t), u_i(t), f_i, dt),
\end{equation}
which numerically integrates $f_i$ by following the dynamics defined in Eq.~\eqref{eq:vel_update}, \eqref{eq:pos_update} and \eqref{eq:cap_update}, respectively. Starting at a state $s_i^{\init} \in \states_i$ and applying a sequence of controls in succession, each for a time step $dt$, gives rise to a dynamically-feasible trajectory $\zeta_i : [0,T] \to \states_i$ where at any $t \in [0,T]$ we have:
\begin{equation}
    \zeta_i(t+dt) \gets \simulate(\zeta_i(t), u_i(t), f_i, dt).
\end{equation}

Let us define a cost function $\mathcal{J}: \states_i \times \mathcal{U}_i \times (\reals_{\geq 0})^m \to \reals_{\geq 0}$ where for a state $s_i(t) \in \states_i$, control $u_i(t) \in \mathcal{U}_i$, and infection levels $\ell(t)$ at time $t$:
\begin{equation}
\mathcal{J}(s_i(t), u_i(t),\ell(t)) = \alpha_1 \|\mu_i(t)\|^2 + \alpha_2 \|v_i(t)\| + \frac{\alpha_3}{n} \sum_{j=1}^m \ell_j(t) + \alpha_4
\end{equation}
where $\alpha_1 \|\mu_i(t)\|^2$ is the energy cost, $\alpha_2 \|v_i(t)\|$ is the movement cost, $\frac{\alpha_3}{n} \sum_{j=1}^m \ell_j(t)$ is the infection cost, $\alpha_4$ is the step cost, and $\alpha_1, \alpha_2, \alpha_3, \alpha_4 \in \reals_{\geq 0}$ are some fixed weights. The objective is to jointly control all robots to eliminate all infection sites while minimizing the costs and avoiding inter-robot collisions. Therefore, we formulate the following stochastic constrained optimization problem:
\begin{flalign}
 & \min_{u_1, \ldots, u_n} \mathbb{E}[ \sum_{t=0}^{T} \sum_{i=1}^{n} \mathcal{J}(\zeta_i(t), u_i(t),  \ell(t)) \mid \zeta_i(0) = s_i^{\init}, \ell(0) = \ell^{\init} ] \label{eq:opt}  \\
\text{s.t.,} \nonumber \\
\text{(i) } & \forall \mathcal{I}_j \in \mathcal{I} : \ell_j(T) < \ell_{\epsilon}, \label{eq:const1} \\
\text{(ii) } & \forall t \in [0,T]: p_i(t) \in \wspace, \quad \text{where,  } \zeta_i(t) = \langle p_i(t), c_i(t), v_i(t) \rangle , \label{eq:const2} \\
\text{(iii) } & \forall t \in [0,T], \forall (\mathcal{M}_i \neq \mathcal{M}_j): \mathcal{A}_i(t) \cap \mathcal{A}_j(t) = \emptyset, \label{eq:const3}
\end{flalign}
where $\ell_{\epsilon} \approx 0$ is an infection-completion threshold, and $\mathcal{A}_i(t), \mathcal{A}_j(t) \subset \wspace$ is the area occupied by robots $i,j$ at time $t$ in $\wspace$. Here, the first constraint ensures that all $m$ infected locations are sprayed. The second constraint ensures that the robots remain within the workspace (agricultural field). And the third constraint ensures that there are no collisions. The definition of $\zeta$ ensures that the robots follow the dynamics.

\section{Environment Design}
\label{sec:mdp}
The optimization problem is formulated as a Partially Observable Markov Decision Process (POMDP) $M = (S, A, \Omega, O, \trans, R, \gamma)$. The \textbf{State Space} $S = \mathcal{S}_1 \times \ldots \times \mathcal{S}_n \times \reals_{\geq 0}^m \subseteq \reals^{(2d+1)n+m}$ captures the joint state $s_t = ( \{ s_i(t) \}_{i=1}^n,\; \{  \ell_j(t) \}_{j=1}^m )$. The \textbf{Action Space} $A = \mathcal{U}_1 \times \ldots \times \mathcal{U}_n$ defines joint actions $a_t = \langle u_1, \dots, u_n \rangle$ where $u_i = \langle \mu_i, \beta_i \rangle$. To model sensor inaccuracies, the \textbf{Observation Space} $\Omega$ provides a noisy estimate $o_t \in \Omega$ of the partially normalized state:
\begin{equation}
    o_t = \Big( \{ (p_i(t), \tilde{c}_i(t), \tilde{v}_i(t)) \}_{i=1}^n,\; \{ \tilde{\ell}_j(t) \}_{j=1}^m \Big) + \mathcal{N}(\mathbf{0}_{(2d+1)n+m}, \sigma_{obs}^2 \mathbf{I}_{(2d+1)n+m}), \label{eq:observation}
\end{equation}
where $\tilde{c}_i(t) = c_i(t) / c_i^{\init}, \tilde{v}_i(t) = v_i(t) / (v_{\max}-v_{\min})$, $\tilde{\ell}_j(t) = \ell_j(t) / \ell_j^{\text{init}}$ and $\sigma_{obs}^2$ is the joint variance of position, capacity, velocity and infection-levels. The \textbf{Observation Function} $O: S \to \Delta(\Omega)$ maps each true state to a distribution over observations, realized by adding Gaussian noise as in Eq.~\eqref{eq:observation}. The \textbf{Transition Function} $\trans: S \times A \to \Delta(S)$ governs the stochastic dynamics, including wind and action noise via Eq. \eqref{eq:vel_update}--\eqref{eq:pos_update}, and infection/capacity updates via Eq. \eqref{eq:iloc_update}--\eqref{eq:cap_update}. The \textbf{Reward Function} $R: S \times A \to \reals$ balances infection reduction against collision avoidance and path efficiency. Let $r_s, r_m, r_L \in \reals_{\geq 0} $ be some small, medium, and very large rewards ($r_s < r_m \ll r_L$). The total reward is:
\begin{alignat}{2}
    R(s_t, a_t) &= \rspray(t) + \rterm(t) + \rpath(t) \label{eq:reward}\\
    \text{where:} \quad \rspray(t) &= r_m \sum_{j=1}^m (\ell_j(t) - \ell_j(t+dt)) \label{eq:r_spr} \\
    \rterm(t) &= 
    \begin{cases}
        -r_L, & \text{if } \exists (\mathcal{M}_i \neq \mathcal{M}_j): \mathcal{A}_i(t) \cap \mathcal{A}_j(t) \neq \emptyset \\
        +r_L, & \text{else if } \sum_{j=1}^m \ell_j(t) < \ell_{\epsilon} \\        
        0, & \text{otherwise}
    \end{cases} \label{eq:r_term} \\
    \rpath(t) &= - \textstyle \sum_{i=1}^n \mathcal{J}(s_i(t), u_i(t), \ell(t)) - r_m \sum_{i=1}^n \mathds{1}_{p_i \notin \wspace} \label{eq:r_path}
\end{alignat}
Finally, the \textbf{Discount Factor} $\gamma \in [0, 1)$ incentivizes rapid infection elimination. We define the following theorem to describe the relation between the optimization problem and the developed POMDP:

\begin{theorem}
\label{thm:optimization}
In the limit $\gamma \to 1$, the optimal policy $\pi^*$ for the POMDP $M$ starting at state $s = (\{ s_i^{\mathrm{init}} \}_{i=1}^n, \ell^{\init} )$  generates trajectories $\{ \zeta_i \}_{i=1}^n$ that solve the optimization problem described in equation (\ref{eq:opt}) with probability close to 1. Conversely, any solution of equation (\ref{eq:opt}) induces an optimal policy $\pi^*$ for the POMDP $M$ with probability close to 1.
\end{theorem} % starting at state $s = (\{ s_i^{\mathrm{init}} \}_{i=1}^n, \{ \ell_j^{\init} \}_{j=1}^m )$

% \begin{theorem}
% \label{thm:optimization}
% The optimal policy $\pi^*$ for the POMDP $M$ will generate the trajectories $\{ \zeta_i \}_{i=1}^n$ starting at state $s = (\{ s_i^{\mathrm{init}} \}_{i=1}^n, \ell^{\init} )$ that solve the optimization problem described in equation (\ref{eq:opt}) with probability close to 1 if $\gamma \approx 1$. Conversely, any solution of equation (\ref{eq:opt}) induces an optimal policy $\pi^*$ for the POMDP $M$ with probability close to 1 if $\gamma \approx 1$.
% \end{theorem} % starting at state $s = (\{ s_i^{\mathrm{init}} \}_{i=1}^n, \{ \ell_j^{\init} \}_{j=1}^m )$

\textbf{Proof.} The proof below conditions on the true state $s_{t}$, which is standard practice when establishing reward-shaping equivalences; it corresponds to analyzing the induced MDP on beliefs in the limit of a sufficient statistic. For a finite time horizon $[0,T]$, the definition of optimal policy $\pi^*$ and optimal action-value function $Q^{\pi^*}$ \citep{sutton2018reinforcement} are as follows:
\begin{align}
& \pi^*(s_t) = \argmax_{a}Q^{\pi^*}(s_t,a) \label{eq:opt_pi}\\
& Q^{\pi^*}(s_t,a_t) = \max_{\pi} Q^{\pi}(s_t,a_t) = \max_{\pi} \mathbb{E}_{\pi} [\sum_{\tau=0}^T \gamma^{\tau} R(s_{\tau},a_{\tau})  \mid s_t, a_t] \ \ \text{[Definition of $Q^{\pi}$ \cite{sutton2018reinforcement}]} \label{eq:opt_policy} \\
& = \max_{\pi} \Big( \mathbb{E}_{\pi} [\sum_{\tau=0}^T \gamma^{\tau} \rspray(\tau) \mid .] + \mathbb{E}_{\pi} [ \sum_{\tau=0}^T \gamma^{\tau} \rterm(\tau) \mid .] + \mathbb{E}_{\pi} [ \sum_{\tau=0}^T \gamma^{\tau} \rpath(\tau) \mid .] \Big) \tag*{[From Eq. \eqref{eq:reward} and properties of $\mathbb{E}$]}  
\end{align}

We prove both directions of the equivalence between the POMDP M and Eq. \eqref{eq:opt} as follows:

\textbf{(Direction $\Rightarrow$) $\pi^*$ solves the optimization problem}: We proceed in two lemmas.

\begin{lemma}
\label{lem:const}
Any optimal policy $\pi^*$ must satisfy constraints (i), (ii), and (iii) of Eq. \eqref{eq:opt} with probability close to one.
% almost surely (probability $\approx 1$).
\end{lemma}

\textit{Proof.} Because the action and position noises ${\epsilon}_a, {\epsilon}_p$ are unbounded, no policy can eliminate the possibility of a constraint violation with probability exactly 1. Instead, we show that the reward shaping makes violations overwhelmingly suboptimal. We analyze each constraint:

\paragraph{Constraint (iii) -- no collisions.} Suppose $\pi^*$ induces a collision at some time $t^*$, i.e., $ \mathcal{A}_i(t^*) \cap \mathcal{A}_j(t^*) \neq \emptyset$ for some $\mathcal{M}_i\neq \mathcal{M}_j$. Then $\rterm = -r_L$. We choose $r_L \gg \sum_{\tau=0}^T \gamma^{\tau} (\rspray(\tau) + \rpath(\tau))$, such that $-r_L$ dominates other rewards. This incurs $Q^\pi \approx - r_L$, which is the lower bound. No such policy can be optimal.

\paragraph{Constraint (ii) -- workspace containment.} Any policy that places a robot $i$ outside $\wspace$ will satisfy the condition $\mathds{1}_{p_i \notin \wspace}$, and have a per-step penalty of $-r_m$ for the robot in the term $- r_m \sum_{i=1}^n \mathds{1}_{p_i \notin \wspace}$ of $\rpath$. Therefore, keeping each robot within $\wspace$ will have $- r_m \sum_{i=1}^n \mathds{1}_{p_i \notin \wspace} = 0$, maximizing $\rpath$ since $r_m > 0$.

\paragraph{Constraint (i) -- infection elimination.} A policy that terminates with $\sum_{j=1}^m \ell_j(T) < \ell_{\epsilon}$ receives $\rterm = +r_L$. Since $r_L \gg \sum_{\tau=0}^T \gamma^{\tau} (\rspray(\tau) + \rpath(\tau))$, achieving $+r_L$ dominates other rewards. Therefore, maximizing $\rterm$ will result in $\ell_j(T) < \ell_{\epsilon} \approx 0$, satisfying constraint (i) by eliminating all infections.

\begin{lemma}
\label{lem:cost}
Among all policies satisfying constraints (i)-(iii) of Eq. \eqref{eq:opt}, maximizing $Q^{\pi}$ is equivalent to achieving the minimization objective of Eq. \eqref{eq:opt}.
\end{lemma}

\textit{Proof.} By Lemma~\ref{lem:const} and constraint~(i), $\rterm(\tau) = +r_L$ at termination and $0$ otherwise; therefore, $\mathbb{E}_{\pi}[\sum_{\tau=0}^T \gamma^{\tau} \rterm(\tau) \mid \cdot] = \mathbb{E}_{\pi}[ \gamma^{T} r_L \mid \cdot] \approx r_L$ where $T$ is the termination time and $\gamma \approx 1$. From Eq.~\eqref{eq:r_spr}:
\begin{flalign}
&\mathbb{E}_{\pi} [\sum_{\tau=0}^T \gamma^{\tau} \rspray(\tau) \mid .] 
= \mathbb{E}_{\pi} [r_m \sum_{\tau=0}^T \gamma^{\tau} \sum_{j=1}^m (\ell_j(\tau) - \ell_j(\tau+dt)) \mid .] \nonumber \\
&\approx \mathbb{E}_{\pi} [r_m \sum_{j=1}^m(\gamma^0\ell_j(0) - \gamma^{T}\ell_j(T)) \mid .] \nonumber = \mathbb{E}_{\pi} [r_m \sum_{j=1}^m(\ell_j^{\mathrm{init}} - 0) \mid .] = r_m L^{\init} \tag*{[Since $\gamma \approx 1$ and $\ell_j(T) \approx 0$]} \nonumber
\end{flalign}
By constraint~(ii), $p_i(\tau) \in \wspace$ for all $i$. Hence $\mathds{1}_{p_i \notin \wspace} = 0$ and from Eq.~\eqref{eq:r_path}:
\begin{flalign}
\mathbb{E}_{\pi}  [\sum_{\tau=0}^T \gamma^{\tau} \rpath(\tau) \mid \cdot] &= -\mathbb{E}_{\pi}  [\sum_{\tau=0}^T \gamma^{\tau} \sum_{i=1}^n (\mathcal{J}(s_i(\tau), u_i(\tau), \ell(\tau))) \mid \cdot]
\end{flalign}
Finally, from Eq.~\eqref{eq:opt_policy}:
\begin{flalign}
Q^{\pi^*}(s_{t},a_{t}) &\approx \max_{\pi} \Big( r_L + r_m L^{\init} - \mathbb{E}_{\pi}  [\sum_{\tau=0}^T \gamma^{\tau} \sum_{i=1}^n \mathcal{J}(s_i(\tau), u_i(\tau), \ell(\tau)) \mid \cdot ] \Big) \nonumber \\
&= r_L + r_m L^{\init} - \min_{\pi} \Big( \mathbb{E}_{\pi}  [\sum_{\tau=0}^T \gamma^{\tau} \sum_{i=1}^n \mathcal{J}(s_i(\tau), u_i(\tau), \ell(\tau)) \mid \cdot] \Big) \tag*{[$\max_a(C-f(a)) = C -\min_a f(a)$]}
\end{flalign}
Here $r_L + r_m L^{\init}$ is approximately constant (policy-independent). Consequently, the remaining degree of freedom in $Q^{\pi^*}$ is precisely the minimization of the expected discounted cost, which proves the lemma for $\gamma \approx 1$ and starting at state $s_t = (\{ s_i^{\mathrm{init}} \}_{i=1}^n, \ell^{\init} )$. Therefore, by Lemmas \ref{lem:const} and \ref{lem:cost}, the optimal policy $\pi^*$ generates the trajectories $\{ \zeta_i \}_{i=1}^n$ that solve the optimization problem described in equation (\ref{eq:opt}).

\textbf{(Direction $\Leftarrow$) the optimization problem induces optimal policy $\pi^*$}: Let $\{ \zeta_i^* \}_{i=1}^n$ be an optimal solution to Eq. \eqref{eq:opt}-\eqref{eq:const3}. Define a time-varying (or trajectory-tracking) policy $\hat{\pi}$ as follows: at each time $t$, given state $s_{t}$, select $\hat{\pi}(s_{t}) = \langle u_1^*(t), ..., u_n^*(t) \rangle$ where $u_i^*(t) = \langle \mu_i^*(t), \beta_i^*(t) \rangle$ is the control along the optimal trajectory $\zeta_i^*$. By Lemma \ref{lem:const}, $\zeta_i^*$ satisfies the three constraints of Eq \eqref{eq:opt} with probability close to one.

\textbf{Claim $\hat{\pi} = \pi^*$}: Since $\zeta_i^*$ satisfies constraint (iii), $\hat{\pi}$ never triggers $-r_L$. Since it satisfies constraint (i), it achieves $+r_L$ at termination. Therefore, $\hat{\pi}$ achieves the maximum rewards for $\rterm$ and $\rspray$. Since it satisfies constraint (ii), $- r_m \sum_{i=1}^n \mathds{1}_{p_i \notin \wspace} = 0$ in $\rpath$. Now suppose for contradiction that there exists another policy $\pi'$ with $Q^{\pi'} > Q^{\hat{\pi}}$. By Lemma~\ref{lem:const} and the definition of $\pi^*$, $\pi'$ must also satisfy all three constraints. By Lemma~\ref{lem:cost}, $Q^{\pi'} > Q^{\hat{\pi}}$ implies:
\[
\sum_{t=0}^T \sum_{i=1}^n \mathcal{J} \Big( \zeta_i^{\pi'}(t), u_i^{\pi'}(t), \ell^{\pi'}(t) \Big)
< \sum_{t=0}^T \sum_{i=1}^n \mathcal{J} \Big(\zeta_i^*(t), u_i^*(t), \ell^*(t) \Big),
\]
contradicting the optimality of $\zeta_i^*$ from Eq. \eqref{eq:opt}. Hence $\hat{\pi} = \pi^*$. Therefore, Theorem 1 holds in both directions. \hfill$\square$

\section{Experiments}
\label{sec:exp}

We present a comprehensive experimental evaluation to measure the following: performance of different RL methods, sensitivity to design choices, out-of-distribution (OOD) robustness, and the sim-to-real gap. The objective is to compare the learning efficiency of different state-of-the-art off-the-shelf RL algorithms across different environmental variations and training conditions.

\paragraph{Setup.}
We compare four on-policy (A2C \citep{mnih2016a2c}, ARS \citep{mania2018ars}, PPO \citep{schulman2017ppo}, TRPO \citep{schulman2015trust}) and two off-policy (CrossQ \citep{bhatt2024crossq}, TQC \citep{kuznetsov2020tqc}) algorithms via Stable-Baselines3 \citep{stable-baselines3}. The main experiments were performed in two settings: (A) using \textit{default hyperparameters} from Stable-Baselines3, and (B) using \textit{tuned hyperparameters} found by Optuna hyperparameter optimization network~\citep{optuna_2019} (details are given in the supplementary). We vary the number of robots $N \in \{2, 3, 4, 5\}$ and crafted 10 variations of the environment by randomizing the field shape, starting positions, and infected locations. Every algorithm--variation pair is trained for $2\times 10^6$ timesteps with 5 seeds $(0,42,123,2024,9999)$. The interquartile mean (IQM)~\citep{agarwal2021deep} is reported as the primary metric. The mean reward threshold for an optimal policy is designed to be +1800 for all configurations. All experiments were conducted on an HPC cluster with different nodes having GPUs: NVIDIA\textsuperscript{\textregistered} GeForce RTX 1080Ti, 2080Ti, 3090, L4, L40s, A4000, A4500, A6000; CPUs: Intel\textsuperscript{\textregistered} Xeon\textsuperscript{\textregistered} E5-2680v3, E5-2683v4, Gold 6130 (16--48 cores); and RAM: 4-16 GB (each node). Simulations were performed on a local desktop with an AMD Ryzen\texttrademark\ Threadripper 1950X (16-core) CPU, 64~GB RAM, and an NVIDIA\textsuperscript{\textregistered} GeForce RTX 2080 GPU. Each experiment took about 1-6 hours to complete on each node. All independent experiments ran on parallel in the HPC cluster.

\begin{table*}[tbp]
\centering
\footnotesize
\caption{\small IQM of episodic reward under default and tuned hyperparameters over 5\,seeds\,$\times$\,10\,variations\,$\times$\,50 episodes. $\dagger$ = best per column.}
\label{tab:main}
\resizebox{\textwidth}{!}{
\begin{tabular}{l|cccc|cccc}
    \toprule
    \multirow{2}{*}{\textbf{Algorithm}} & \multicolumn{4}{c|}{\textbf{Default Hyperparameters}} 
    & \multicolumn{4}{c}{\textbf{Tuned Hyperparameters}} \\
    \cmidrule(lr){2-9}
           & $N=2$                     & $N=3$                     & $N=4$                     & $N=5$                     
           & $N=2$                     & $N=3$                     & $N=4$                     & $N=5$                     \\
    \midrule    
    A2C    & $-694.4$ & $-573.8$ & $-529.9$ & $-536.7$ & $-669.6$ & $-488.8$ & $-483.6$ & $-489.2$ \\
    ARS    & $-1217.0$ & $-1034.2$ & $-943.2$ & $-852.8$ & $-1726.9$ & $-1254.9$ & $-1182.4$ & $-1053.5$ \\
    PPO    & $-1099.7$ & $370.1$ & $419.1$ & $32.3$ & $-681.0$ & $-738.2$ & $-682.4$ & $-676.8$ \\
    TQC    & $-708.4$ & $-89.7$ & $889.7$ & $1211.8$ & $-9.0$ & $1802.8$ & $\mathbf{1914.7}^\dagger$ & $\mathbf{1931.3}^\dagger$ \\
    TRPO   & $13.1$ & $1799.2$ & $1830.9$ & $1851.0$ & $-533.4$ & $682.5$ & $648.1$ & $111.5$ \\
    CrossQ & $\mathbf{1802.5}^\dagger$ & $\mathbf{1868.1}^\dagger$ & $\mathbf{1896.8}^\dagger$ & $\mathbf{1901.8}^\dagger$ & $\mathbf{1827.2}^\dagger$ & $\mathbf{1895.8}^\dagger$ & $-460.1$ & $-623.4$ \\
    \bottomrule
\end{tabular}
\vspace{-12pt}}
\end{table*}

% ─── Table: observation gap ──────────────────────────────────────────────
% \begin{table}[t]
% \centering
% \caption{Sim-to-real observation-gap analysis (CrossQ + full DR, $N\!=\!3$,
%   variation~1, 50 eval eps from CoppeliaSim). $\Delta$IQM: relative change vs.\ no-perturbation baseline.}
% \label{tab:obs_gap}
% \small
% \begin{tabular}{lcc}
% \toprule
% \textbf{Perturbation} & IQM & $\Delta$\textbf{IQM (\%)} \\
% \midrule
% No perturbation (baseline)     & $624.45$ & $0.0$ \\
% GPS noise ($\sigma\!=\!0.05$\,m)    & $468.92$ & $-24.9$ \\
% Wind estimate latency (1 step) & $-282.11$ & $-145.2$ \\
% Comm.\ packet loss (5\%)       & $367.00$ & $-41.2$ \\
% All combined                   & $-1267.71$ & $-303.0$ \\
% \bottomrule
% \end{tabular}
% \end{table}
% ─────────────────────────────────────────────────────────────────────────────
\subsection{Main Results}
\label{sec:exp:main}
% ─────────────────────────────────────────────────────────────────────────────

Tables~\ref{tab:main} show the comparison of all six algorithms with default and Optuna-tuned hyperparameters by running 50 episodes for evaluation using the trained models. CrossQ consistently achieved the highest IQM across all~$N$ with default HPs and for $N < 4$ with tuned HPs. TQC performs best for $N > 3$ with tuned HPs, confirming that distributional critics benefit from tuning under stochastic noises. TRPO achieves the second-best performance for $N > 2$ with default HPs and for $N > 3$ with tuned HPs. TRPO's advantage over PPO (counterintuitive given standard benchmarks \citep{schulman2017ppo}) is consistent with its hard KL constraint preventing premature convergence in the high-dimensional coordination space~\citep{kuba2022happo, yu2022mappo}. ARS underperforms throughout, consistent with its known difficulty in high-dimensional continuous action spaces~\citep{mania2018ars}. Learning curves from training (Fig.~\ref{fig:curves}) show the superior performance of CrossQ and TQC, where the dotted lines represent the mean reward threshold. 

% ─── Figures (one 3-panel figure replaces 5 separate ones) ───────────────
\begin{figure}[tbp]
    \centering
    \includegraphics[width=\linewidth]{figures/combined_learning_curves.png}
    \caption{Learning curves - Default (top row) vs tuned (bottom row) hyperparameters. The dotted lines are the reward threshold for the optimal policy.}
    \label{fig:curves}
    \vspace{-12pt}
\end{figure}

% ─────────────────────────────────────────────────────────────────────────────
\subsection{Ablation Study}
\label{sec:exp:ablation}
% ─────────────────────────────────────────────────────────────────────────────

% \paragraph{Ablation study.}
We ablate reward terms, observation spaces, and uncertainty parameters using CrossQ ($N\!=\!3$, 5 seeds, 10 variations, 50 episodes). For reward ablations, we compare the IQM of episode lengths where the episodes are terminated upon at least 95\% infection coverage or on collision. Table~\ref{tab:ablations_and_dr} shows that removing $R_{\mathrm{spr}}$ or $R_{\mathrm{term}}$ causes collisions after training (no learning). Omitting $R_{\mathrm{path}}$ has minimal impact, as the discount factor $\gamma$ automatically minimizes path length at some scale. For observation ablations, removing spatial data is catastrophic, and a positions-only configuration drastically reduces IQM, though dropping infection-level information has negligible effect. Regarding uncertainty, adding only actuation noise yields the highest IQM, slightly outperforming deterministic training (no noise). Conversely, wind noise drops performance to $1694.5$, matching full noise ($1714.5$) and confirming wind as the primary source of performance degradation.

% ─────────────────────────────────────────────────────────────────────────────
\subsection{Domain Randomization}
\label{sec:exp:dr}
% ─────────────────────────────────────────────────────────────────────────────

% \paragraph{Domain randomization.}
We implement domain randomization (DR) \citep{tobin2017domain, peng2018sim2real} by re-sampling physical parameters per training episode under three conditions: (A) No DR, (B) Wind DR only, and (C) Full DR (wind, actuation, spray radius, robot mass/thrust). Details of the randomized parameters and bounds are given in the supplementary. We evaluate these conditions on in-distribution ($\text{wind} \in [0, 0.5]$\,m/s) and out-of-distribution (OOD) ($\text{wind} \in (0.5, 2.0]$\,m/s) sets (Table~\ref{tab:ablations_and_dr}). Both DR conditions significantly improve both in-distribution and OOD performance compared to no DR, and wind DR achieves the best IQM. However, full DR minimizes worst-case tail risk, achieving the lowest CVaR$_{0.1}$ \citep{tamar2015cvar} ($\mathbf{-2428.7}^\dagger$), followed closely by wind DR ($-2937.0$). Ultimately, wind is the primary generalization bottleneck; randomizing additional parameters yields diminishing mean OOD returns but maximizes worst-case robustness.
\vspace{-8pt}

% ─────────────────────────────────────────────────────────────────────────────
\subsection{Simulation and Sim-to-real Analysis}
\label{sec:exp:s2r}
% ─────────────────────────────────────────────────────────────────────────────

% ============================================================
% OPTION 1: wraptable  (requires \usepackage{wrapfig})
% Table sits to the RIGHT of the text, text wraps around it.
% Change {r} to {l} to place the table on the left instead.
% The second argument (5.8cm) is the total width of the table
% column — tune this to fit your content.
% ============================================================

\begin{table*}[tbp]
\centering
\small
\caption{\small Ablation study and domain randomization results for $N=3$ using CrossQ, over 5\,seeds\,$\times$\,10\,variations\,$\times$\,50 episodes. Left three blocks: ablation experiments. Right block: DR experiments. $\dagger$ = best per block/column.}
\label{tab:ablations_and_dr}
\resizebox{\textwidth}{!}{%
\begin{tabular}{lcc|lc|lc|lccc}
    \toprule
    \multicolumn{3}{c|}{\textbf{Reward ablation}}
    & \multicolumn{2}{c|}{\textbf{Observation ablation}}
    & \multicolumn{2}{c|}{\textbf{Uncertainty}} & \multicolumn{4}{c}{\textbf{Domain Randomization}} \\
    \midrule
    % \cmidrule(lr){1-3}\cmidrule(lr){4-5}\cmidrule(lr){6-7}
    Condition &  \makecell{Ep. length \\ IQM} & \makecell{Terminal\\ Condition}
    & Condition & \makecell{Reward \\ IQM}
    & Condition & \makecell{Reward \\ IQM}
    & Condition & \makecell{In-dist. \\ IQM} & \makecell{OOD \\ IQM} & CVaR$_{0.1}$ \\
    \midrule
    Full reward & $\mathbf{45.9}^\dagger$ & Sprayed (97\%) & Full & $\mathbf{1881.3}^\dagger$ & Full & $1714.5$ & No DR & $-3302.9$ & $-62156.9$ & $-88938.7$ \\
    No $R_\text{term}$ & $23.5$ & Collision & No pos. & $-973.0$ & Wind only & $1694.5$ & Wind DR & $\mathbf{-601.0}^\dagger$ & $\mathbf{-601.0}^\dagger$ & $-2937.0$ \\
    No $R_\text{spr}$ & $346.2$ & Collision & No infect. & $1865.3$ & Act. only & $\mathbf{1907.7}^\dagger$ & Full DR & $-620.3$ & $-620.3$ & $\mathbf{-2428.7}^\dagger$ \\
    No $R_\text{path}$ & $50.7$ & Sprayed (96\%) & Pos. only & $1408.2$ & Deterministic & $1892.7$ &  &  &  &  \\
    \bottomrule
\end{tabular}}
\vspace{-12pt}
\end{table*}

\begin{wraptable}{r}{7cm}           % {r|l} = side,  5.8cm = column width
\vspace{-1.8em}                       % pull table up to align with paragraph top
\centering
\caption{\small Sim-to-real observation-gap analysis (CrossQ, $N\!=\!3$, variation~1, 50 eval eps).
  $\Delta$IQM: relative change vs.\ no-perturbation baseline.}
\label{tab:obs_gap}
\setlength{\tabcolsep}{4pt}           % tighten column spacing for narrow width
\small
\begin{tabular}{lcc}
    \toprule
    \textbf{Perturbation} & \textbf{IQM} & $\Delta$\textbf{IQM (\%)} \\
    \midrule
    None (baseline)               & $1887.92$   & $0.0$    \\
    GPS noise ($\sigma\!=\!0.05$\,m)   & $1565.35$   & $-17.1$  \\
    Wind latency (3 steps)         & $665.22$  & $-64.8$ \\
    Comm.\ loss (10\%)             & $1853.52$   & $-1.8$  \\
    All combined                  & $-2144.47$ & $-213.6$ \\
    \bottomrule
\end{tabular}
\vspace{-0.8em}                       % reclaim space below table
\end{wraptable}

We deploy trained policies in a 3D CoppeliaSim~\citep{coppeliaSim2013} agricultural scene with PyBullet~\citep{coumans2021pybullet} physics connected via ZeroMQ~\citep{ZeroMQ}. Table~\ref{tab:obs_gap} quantifies IQM degradation under inference-time sensor imperfections applied to the CrossQ policy trained with default hyperparameters ($N\!=\!3$, variation~1). Wind estimate latency (3 steps) is the dominant transfer risk ($\Delta$IQM~$=\!-64.8\%$, IQM $665.22$), reflecting strong policy dependence on timely wind information. GPS position noise has the smallest individual effect ($-17.1\%$), as the spray-coverage reward is spatially smooth. Communication packet loss (10\%) causes only marginal degradation ($-1.8\%$), yet applying all perturbations simultaneously collapses performance to IQM~$-2144.47$ ($\Delta$IQM~$=\!-213.6\%$), consistent with transfer gaps reported for UAV sim-to-real tasks in existing works~\citep{peng2018sim2real, akkaya2019solving}. Fig.~\ref{fig:sim} illustrates the optimal spraying paths for environment variations 1 and 3.

\begin{figure}[tbp]
    \centering
    \subfloat[Env. 1 initial]{\includegraphics[width=0.247\textwidth, height=0.17\textwidth]{figures/simulations/env1_init.png}\label{fig:env1_init}} 
    \hfil
    \subfloat[Env. 1 final]{\includegraphics[width=0.247\textwidth, height=0.17\textwidth]{figures/simulations/env1_v1_best.png}\label{fig:env1_best}} 
    \hfil
    \subfloat[Env. 3 initial]{\includegraphics[width=0.247\textwidth, height=0.17\textwidth]{figures/simulations/env3_v1_init.png}\label{fig:env3_init}}
    \hfil
    \subfloat[Env. 3 final]{\includegraphics[width=0.247\textwidth, height=0.17\textwidth]{figures/simulations/env3_v1_best.png}\label{fig:env3_best}} 
    \caption{\small Simulation of environment variation 1 (a, b) and 3 (c, d) for $N=3$ (three UAVs). The dark green plants denote the infected locations, and the red lines denote the spraying paths.}
    \label{fig:sim} 
    \vspace{-12pt}
\end{figure}

\section{Conclusion}
\label{sec:con}
We explored how to construct RL environments with correctness guarantees for multi-robot selective spraying. Our work provides researchers with a robust foundation and clear framework to develop environments for such tasks. Looking forward, our future plan includes extending this environment for implementing it on different real robots working in an actual agricultural field, and adding more constraints, such as failure or disconnection of some robots, while finding the optimal path for the remaining active robots.


\bibliographystyle{plainnat}
\bibliography{ref}
%%%%%%%%%%%%%%%%%%%%%%%%%%%%%%%%%%%%%%%%%%%%%%%%%%%%%%%%%%%%

\end{document}